In the introduction we have highlighted the key concepts presented in the previous publications on conditional analyses. In this section we describe most relevant to our work ideas from those papers. The earlier research has introduced two ``conditional'' notions -- {\em conditional soundness} and {\em conditional analysis}. In their essence both concepts are the same -- the result of an analysis holds only under the established conditions. They diverge in their approaches in obtaining unconditional results. Conditional soundness relies on another analysis to prove that the program states that satisfy the negated condition are not reachable, i.e., showing the infeasibility of the program execution containing those states. While conditional analysis may use the same or a different analysis to show that those states do not violate properties either by appearing in behaviors that conform to the properties or by belonging to  infeasible program executions.

The work on conditional soundness~\cite{Conway:SAS08} explores the dependency of the soundness of one analysis on the result of another analysis. Many analyses assume that a program is ``well-behaved'', for example an analysis with the abstract domain $\{-,0,+\}$ assumes the absence of overflow conditions when the rule $\{+\} + 1$ produces not expected $\{+\}$ but $\{-\}$. The authors demonstrate this concept on a pointer analysis for C program. A points-to analysis can produce sound results if the program is memory safe, i.e., there is no out of bound accesses. The authors argue that it is possible to combine points-to and memory safe analyses into a single analysis but such approach is not scalable. Thus they propose to first perform points-to analysis assuming that the program is memory safe and after that run the memory safety analysis with the points-to analysis result to prove that the assumption holds.

Besides explaining the concept of conditional soundness the authors also formalize it through a set of definitions. They specify of a conditional soundness with respect to a predicate $\theta$ on a concrete state. Then an analysis is defined to be $\theta$ sound when it over-approximates program behavior reachable through concrete states satisfying $\theta$. Thus an analysis applies its transfer function only to the states satisfying $\theta$. However, the transfer function can produce non-$\theta$ states but those states will not be explored on the next application of the transfer function.  In order to prove an unconditional soundness of the first analysis the second analysis explores the states produced by the first analysis and proves that resulted non-$\theta$ states are unreachable in the program.

The work~\cite{Naik:POPL07} on static race detection for multi-threaded programs is conceptually similar to the previously described paper on conditional soundness. A static race detection algorithm analyzes memory accesses $m_1$ and $m_2$ guarded by locks $l_1$ and $l_2$ respectively. It assumes that $m_1$ and $m_2$ may alias and thus the algorithm aims to prove that $l_1$ and $l_2$ must alias in order for a program to be race free. In the alternative approach proposed in this paper it is assumed that $l_1$ and $l_2$ must not alias and then it is shown that $m_1$ and $m_2$ cannot refer to the same memory location. The absence of aliasing between $m_1$ and $m_2$ is easy to show when $m_i$ is the field of the object referred by $l_i$ and that the field $m_i$ is only reachable through that object. However, in order to make the result of the analysis sound unconditionally another analysis must be prove that $l_1$ and $l_2$ in fact must not alias.

%The main goal of the conditional model checking approach is not to waste resources when a model checker cannot verify the whole program either because of its inability to handle assertion, i.e., not in that domain or running out of resources.
The authors of the paper on conditional model checking~\cite{Beyer:Tech11} emphasize the fact that when a model checker fails either to verify a program or to produce a counterexample then the effort put into the exploration of the program's state space is wasted. These situations arise when either the model checker runs out of computational resources or its unable to decide on the predicate of a branch condition. The latter takes place when the model checker is not designed to handle the theory used in the predicate, e.g., the predicate analysis cannot handle non-linear arithmetic.

In the case of resources exhaustion, i.e., either running out of memory or timing out,  instead of conventional termination the authors propose an informative termination where upon shutting down the model checker produces data that describes the parts of the state space it was able to verify. To characterize the parts of the analyzed state space the model checker enhances each abstract state with a predicate that uniquely describes that state. Hence, when the model checker depletes its computational resources, instead of terminating with a run-time exception, it returns a condition describing the set of states it has analyzed. The condition is expressed as a conjunction of the analyzed states predicates. The next model checker can use the negated condition to guide itself to the part of the state space that was not analyzed by the previous model checker. This characteristic of conditional model checker corresponds to the notion of conditional analysis.

In the case of undecidable predicates the authors propose to proceed with such predicate interpretation that permits the model checker to proceed with the verification, e.g., assuming that an assertion it cannot reason about does hold. The encountered assumptions for taken branches are encoded as an automaton that is used to guide the further analysis to the assumed branches. These assumptions must be analyzed further by another model checker that is capable to reason about them. If the subsequent model checker determines that these assumptions hold then the verification is complete. This side of conditional model checker relates to the idea of conditional soundness.

%Moreover, this conditional approach allows a model checker to express its bounds like the tree depth, loop unrolling bounds and other boundary criteria as a condition. Thus if a model checker was not able to find  Also when a model checker is bounded by the depth of a path or location repetition the technique is able to translate these bounds into conditions which can describe for what states the model checker was able to verify the program.

A conditional model checker integrates assumptions, i.e., conditional soundness for certain states, and predicate , i.e., conditional analysis, information into the characteristic of abstract states. At the end of the execution  the conditional model checker traverses the reached abstract states and produces the condition and the branch outcome assumption information. The experiments presented in the paper have shown that applying conditional model checkers with feedback capabilities increased the number of verified programs.

%The fundamental difference with our approach is model checking is MOP analysis where each path is analyzed separately. Thus if an analysis cannot decide something is not because it lost information on a merge point, the analysis cannot handle it inherently like a non-linear arithmetic. In MOP one can specify additional bounds as loop unrolling, resource time out, depth of a search tree to be quantified as a condition.

%In MFP analysis one cannot have condition as loop unrolling (inapplicable), or timeout (not sound). In MFP imprecision may arise due to merging operator on control flows.

Our work differs from the previous research in the approach of determining the condition for which an analysis guarantees its most accurate result. In the work on conditional soundness these conditions are defined a priori and applied for all states, e.g., no integer overflows. These conditions are derived straight from the design of the analysis. Without those conditions in place the analysis is unable to produce a correct answer. A conditional model checker discovers the conditions, i.e., the assumptions, for its optimal result while analyzing a program. Although its conditions only applicable to some set of states, e.g., the false branch of the conditional statement on line 10 is infeasible. 

However, conditional data-flow analysis cannot use neither of these approaches. Clearly, finding the conditions solely from the design of data-flow analysis is not possible because these conditions are program dependent. Therefore no a priori conditions can be determined. Determining them during analysis run as done in the conditional model checking is also problematic. The reason is that model checking and data-flow analysis have different underlying analysis frameworks. Model checking is based on meet over path (MOP) analysis framework where an analysis calculates facts for each program path separately. While data flow analysis is based on maximal fixed point (MFP) analysis framework where only a single fact is calculated for a set of paths and that fact over-approximates the facts of the evaluated paths.

Thus, when an MOP analysis cannot decide on a branch outcome it is solely due to the analysis theory being inadequate. However, for an MFP analysis an imprecision may also be the result of the over-approximating several paths with diverse facts. But the MFP analysis might produce a better result if the set of paths it over-approximates is smaller with more uniform facts. The current work looks into exploring alternative methods in partitioning program paths to increase the precision of data flow analysis.  

%In our case $\theta$ can be describes the set of states reachable from the restricted input domain or lay on the selected set of paths. In our case it is conditional analysis not conditional soundness.

%Also it combines different conditions as assumptions, location monitoring, overflow monitoring. For assumption the condition is an expression, for other the condition for an infeasible path is ``false'' which exclude that path. Also track predicates for the case of resources exhausting which tells what paths have been covered.

%Similar to conditional model checking, conditional soundness gives the condition under which analyses are most precise. Without this condition the analysis cannot produce any result, i.e., always unknown, and with that condition the analysis is able to produce some result. Soundness here is more related to the precision. Once again this is different from our work where we do not know when DF are the most precise.

%Similar that the result produces by conditional and data-flow analysis are only sound within that condition. 

%In order to prove unconditional soundness some other analysis must prove that non-$\theta$ states are not feasible. In our work for unconditional soundness we need to prove that non-$\theta$ states do not violate properties which also include showing that some non-$\theta$ states are not feasible.

%Work on the conditional model checking the unconditional soundness of the program is achieved as a hybrid. For assumption it is to show that non-$\theta_{assumption}$ states are not feasible and for partial result is to show that non-$\theta_{predicate}$ do not violate that condition. We just need to show the latter case. 















